\subsection{Feedback tests for owner closures}
\label{app:core-feedback}

Pure-owner closures satisfy a closed feedback-elimination law.
With two foreign owners its repeated queries compress to a single
mutual-capture test. With three foreign owners this compression fails;
an eight-variable example detects the missing unused-owner test.
Even after every such extension, a five-variable terminal tuple passes
feedback but has an unrealizable pure source port.

For a core operation on essential support $V$, its pure-owner closure
$c_r$ can be recovered from a one-owner trace. Replace foreign gates
by union and distribute to obtain a sum of $r$-folds on nonempty sets
$\mathcal H_r$. This preserves the pure-$r$ source coefficients and
the $r$-profile, and
\[
 c_r(E)=V\setminus\bigcup\{H\in\mathcal H_r:H\cap E=\varnothing\}.
\]
The union on the right is an interior on complements, so $c_r$ is a
loopless closure. Its spanning predicate is the singleton profile:
$\Phi_r(E)=[c_r(E)=V]$.

Here is the feedback law for arbitrary loopless closure data. Fix an
observing owner $r$, a coalition $A$, and finitely many queries
$(s_i,X_i,P_i)$ with $s_i\ne r$, such that
\begin{align*}
 &X_i\cap P_i=\varnothing,\qquad X_i\subseteq c_{s_i}(P_i),\\
 &\bigcup_i(X_i\cup P_i)\subseteq c_r(A).
\end{align*}
For each $i$, require $A\cap P_i$ to be covered by some of the $X_j$
having pairwise distinct owners, all different from $s_i$. Then impose
\begin{equation}\label{core:eq:feedback-elimination}
 c_r(A)=c_r\left(A\setminus\bigcup_iX_i\right).\tag{FE}
\end{equation}
All quantifiers in this law range over the supplied finite closure data.

Call $B$ \emph{$r$-independent} if
$b\notin c_r(B\setminus\{b\})$ for every $b\in B$. Put
$U=c_r(B)$, $Z=U\setminus B$, and define the conditional foreign closures
\begin{equation}\label{core:eq:conditional-foreign-closure}
 d_s(E)=c_s(Z\cup E)\cap B\qquad(s\ne r,\ E\subseteq B).
\end{equation}
These are closures even when $U$ is not foreign-closed: if
$G=c_s(Z\cup E)$, then $Z\cup d_s(E)\subseteq G$, so
$d_s(d_s(E))\subseteq G\cap B=d_s(E)$; extensivity gives equality.

An equivalent independent-base form of (FE) forbids the following
\emph{feedback family}: sets $X_i,Q_i\subseteq B$ with
\begin{equation}\label{core:eq:base-feedback}
 X_i\cap Q_i=\varnothing,\qquad X_i\subseteq d_{s_i}(Q_i),
 \qquad \bigcup_iX_i\ne\varnothing,
\end{equation}
where each $Q_i$ is covered by head blocks $X_j$ of pairwise distinct
owners different from $s_i$. Indeed a forbidden family lifts to (FE)
using $A=B$ and $P_i=Z\cup Q_i$, contradicting independence.
Conversely, for a failing (FE) instance choose a minimal
$B\subseteq A$ with $c_r(B)=c_r(A)$. It is $r$-independent.
Replace $X_i,P_i$ by $X_i\cap B,P_i\cap B$ and use
$Z=c_r(B)\setminus B$. Capture and the covers restrict to $B$;
some restricted head is nonempty, since otherwise
$B\subseteq A\setminus\bigcup_iX_i$ would prove (FE).

\begin{proposition}[Necessary feedback elimination]
\label{core:prop:feedback-elimination}
The pure-owner closures of every core operation satisfy (FE).
\end{proposition}
\begin{proof}
Use its independent-base form. The neutral residual on $U=c_r(B)$
is core by \cref{core:thm:neutral-elimination}, its closures are
$c_s(E)\cap U$, and its $r$-profile is positive at $B$.
At query $i$, give $X_i$ the menu $I$, give $Z\cup Q_i$ the $s_i$
atom, and every other input $I\setminus\{r\}$. Its only possible
$r$-sources lie in $X_i$, whose pure-$s_i$ ports vanish; hence the
residual omits $r$ there.

An omission budget with head $X$ and owner $s\ne r$ forces the
positive profile $\Phi_s(U\setminus X)$: decrease every input outside
$X$ to $\{s\}$. Monotonicity preserves omission of $r$, and the new
output lies in $\{\{s\},I\}$, so must be $\{s\}$.

Now pass the $B$-positive $r$-profile and all budgets down the word.
At an $r$-root choose a positive child; all budgets pass to both.
At a foreign root $s$, both children are $B$-positive, and budgets
of other owners pass directly. A child's $s$-profile cannot be
positive on $Z\cup Q_i$ for an $s$-colored query: together with
its $r$-profile at $B$ and the forced profiles of the rainbow budgets
covering $Q_i$, the common-source condition would have empty
intersection. Thus neither child absorbs at that query, and its
budget passes too. Induction gives an $r$-positive profile on
$B\setminus\bigcup_iX_i$: a positive leaf lies in $B$, and its budgets
exclude every head block. Consequently this set spans $U$, contrary
to independence and the nonempty head union.
\end{proof}

For closures $d_s$ on a finite set $B$, the \emph{hereditary profile
test} requires that, for every nonempty $L\subseteq B$, spanning
coalitions $A_s\subseteq L$ with $L\subseteq d_s(A_s)$ have nonempty
common intersection. In the conditional setting, (FE) implies this
test: use one query per owner with $X_s=L\setminus A_s$, $Q_s=A_s$.
If the intersection were empty, each premise would have a rainbow
cover by the other heads.

\begin{theorem}[Two-color feedback compression]
\label{core:thm:two-color-feedback}
For two closures $d_1,d_2$ on a finite set, the hereditary profile
test is equivalent to absence of every finite feedback family. It is
also equivalent to
\begin{equation}\label{core:eq:conditional-mutual-capture}
 C\cap D=\varnothing,\quad C\subseteq d_1(D),\quad
 D\subseteq d_2(C)\quad\Longrightarrow\quad C=D=\varnothing.
\end{equation}
In particular, for $K=3$, (FE) is equivalent to the hereditary profile
test at every observer-independent base.
\end{theorem}
\begin{proof}
Both the profile test and the mutual-capture condition exclude loops,
using a singleton retained set and allowing one side to be empty.
For loopless closures the latter is exactly the hereditary profile
test: a disjoint spanning pair gives mutual capture, and a nonempty
mutual pair spans its own union. Such a pair is itself a two-query
feedback family.

Suppose the profile test holds but a feedback family exists.
A nonempty head has a nonempty premise, covered by one opposite-color
head. Following covers gives an alternating directed cycle. Replace
each cycle head by the preceding query's premise to obtain nonempty
sets, indexed cyclically modulo $m$, with
\[
 A_i\subseteq d_1(B_i),\quad B_i\subseteq d_2(A_{i+1}),\quad
 A_i\cap B_i=B_i\cap A_{i+1}=\varnothing.
\]
Put
\[
 C=\bigcup_{i=0}^{m-1}\left(B_i\setminus\bigcup_{h<i}A_h\right),
 \qquad
 D=\bigcup_{i=0}^{m-1}\left(A_{i+1}\setminus\bigcup_{h<i}B_h\right).
\]
Starting with $C$, closure $d_1$ successively generates $B_0,A_0$,
then $B_1,A_1$, and so on; deleted elements are already generated.
Starting with $D$, closure $d_2$ similarly generates $A_1,B_0$,
then $A_2,B_1$, and so on. Thus both closures contain the full cycle
union, giving $D\subseteq d_1(C)$ and $C\subseteq d_2(D)$.
Both sets are nonempty.
If an element belonged to their $i$- and $j$-summands respectively,
the exclusions would force $i\geq j$ and $(j+1\bmod m)\geq i$.
Without wraparound this gives $i=j$ or $i=j+1$, both excluded by
adjacent disjointness; wraparound is impossible except $m=1$, also
excluded by disjointness. Hence $(D,C)$ is a forbidden mutual pair.
The last assertion follows from the independent-base form of (FE).
\end{proof}

\begin{theorem}[A three-color feedback obstruction]
\label{core:thm:three-color-feedback}
There is an eight-variable loopless closure triple satisfying (FE)
but not jointly realizable as the pure-owner closures of a three-owner
core operation. After adjoining an identity owner, every conditional
hereditary profile test still holds, while (FE) fails on six queries.
Under any extension by identity owners the tuple also satisfies the
simultaneous cross-value composition law
\begin{equation}\label{core:eq:feedback-cross-composition}
 P\subseteq c_r(E)\setminus E\quad\Longrightarrow\quad
 c_s(P)\subseteq c_r(E).
\end{equation}
\end{theorem}
\begin{proof}
Use owners $a,b,c$ and variables
$u_a,u_b,u_c,p,q,e,t_+,t_-$. Each closure is generated by its two rules:
\[
\begin{array}{c|c|c}
\text{query}&\text{premise }Q_i&\text{head }X_i\\\hline
 a_0&u_b,u_c,q,t_-&u_a,p,t_+\\
 a_1&u_b,u_c,p,t_+&u_a,q,t_-\\
 b_0&u_a,u_c,e,t_+&u_b,p,t_-\\
 b_1&u_a,u_c,p,t_-&u_b,e,t_+\\
 c_0&u_a,u_b,e,t_-&u_c,q,t_+\\
 c_1&u_a,u_b,q,t_+&u_c,e,t_-
\end{array}
\]
The two queries of each color share the same set $Q_i\cup X_i$.
Once either rule fires, both do; writing $H_s=X_{s_0}\cup X_{s_1}$,
the closure is $E$ if inactive and $E\cup H_s$ otherwise.
A first firing adds at most its three head variables.

First verify the hereditary profile test for this triple. If some
colors are active at their spanning coalitions and some are not,
$u_t$ for an inactive color $t$ lies in every active premise;
inactive coalitions equal the retained set. This gives a common
variable. The all-inactive case is immediate.
For the all-active case, regard
$a_0=b\to a,a_1=c\to a,b_0=a\to b,b_1=c\to b,c_0=a\to c,c_1=b\to c$.
The variables $u_s$ are incoming stars; $p,q,e$ are directed two-cycles
and $t_+,t_-$ are the directed triangles. A head contains exactly the
configurations using its arc; $Q_i$ is the union of its owner's two
outgoing-arc heads with $X_i$ removed. Delete one first-fired incoming arc per
owner. The remaining graph has one incoming arc per vertex and hence
a directed cycle. Its configuration belongs to the premise of every
color on that cycle, and is outside $H_s$ for a color off it, so spanning
forces it into that coalition too. It is a common variable.

Now verify every observer-independent-base test, including after
adding identity owners. An identity observer, or an inactive independent
base, has $Z=\varnothing$; the preceding test applies, taking the
omitted color's coalition to be the entire retained set.
An active independent base is $Q_i$ or $Q_i$ plus the opposite
two-cycle variable: any other extra variable is already captured from
$Q_i$ and would be redundant. Color symmetry reduces to $i=a_0$.
Writing $b'=u_b,c'=u_c,v=t_-$, the bases are
$B_- =\{b',c',q,v\}$ and $B_+=B_-\cup\{e\}$, with
$Z=\{u_a,p,t_+\}$. Their two nonidentity conditional closures have the
following rules, with juxtaposition denoting a set:
\[
\begin{array}{c|l|l}
 &d_b&d_c\\\hline
 B_-&c'v\Rightarrow b'&b'q\Rightarrow c'v\\
 B_+&\begin{gathered}c'e\Rightarrow b'v\\c'v\Rightarrow b'e\end{gathered}
     &\begin{gathered}b'ev\Rightarrow c'q\\b'q\Rightarrow c'ev\end{gathered}
\end{array}
\]
For $B_-$, $q$ is uncaptured by both; if it is absent from a retained
set, $d_c$ is identity there. For $B_+$, two disjoint spanning
coalitions would both activate. The $d_c$ coalition cannot start with
$b'q$, since $q$ is uncaptured by $d_b$; hence it contains $e,v$.
But $d_b$ needs one of $e,v$ to activate, a contradiction. All these
closures are loopless, and extra identity coordinates have no effect.
This proves the asserted conditional tests. For three owners,
\cref{core:thm:two-color-feedback} therefore gives (FE).

Add an identity observer $\rho$. At its independent base $V$, use all
six displayed queries. Each premise is covered by the two heads of
arcs leaving its owner, whose colors are distinct and foreign; their
head union is $V$. Thus \eqref{core:eq:base-feedback} violates (FE).
If a three-owner word realized the original triple, interpreting the
same word with fresh owner $\rho$ would preserve all old one-owner
traces. Its $\rho$-trace is the union of every variable, since its
support is $V$ and no gate is owned by $\rho$, giving identity closure.
This would realize the forbidden extension, contradicting
\cref{core:prop:feedback-elimination}. The triple is not realizable.
Finally a set $P$ in \eqref{core:eq:feedback-cross-composition} has
at most three elements. It contains no four-element rule premise,
so every $c_s$ fixes it, proving that law.
\end{proof}

Thus repeated-query aggregation is determined by conditional profiles
for two foreign colors and requires additional information for three.
Pure-port compatibility must also survive adjoining an identity owner.
One fresh identity owner already tests every extension by unused owners:
an identity-colored query has $X_i\subseteq P_i$ and hence empty head,
and all identity observers impose the same remaining conditions.
These statements concern the displayed closure tests; full source-guard
compatibility requires more even when every identity extension passes them.

\begin{theorem}[Identity-stable feedback is insufficient]
\label{core:thm:identity-stable-feedback}
For every finite owner set containing three distinct owners there is a
five-variable loopless closure tuple satisfying (FE) after every extension
by identity owners, but having no joint pure-port realizer. Its separable
reconstruction has full colored-cofactor independence and core proper
nonempty neutral residuals.
\end{theorem}
\begin{proof}
Use the tuple of \cref{core:thm:terminal-separable-obstruction}, with
old support $W=\{x,y,z,t\}$ and new source $X$. Nonrealizability and
all the reconstruction claims were proved there; it remains to check (FE).

Fix an observer $r$ and an $r$-independent base $B$, and put
$U=c_r(B)$ and $Z=U\setminus B$. If $U=\varnothing$, no feedback
family can have a nonempty head union. If $U$ is nonempty and proper,
the tuple restricted to $U$ comes from a core proper
neutral residual. Its conditional closures are exactly those for the
original base, so it passes the feedback test. We need only the
independent bases with full closure.

For observer $a$ these are $\{t,X\}$ and $J\cup\{X\}$, where
$J$ is a two-element subset of $\{x,y,z\}$. For observer $b$ they
are $J\cup\{t,X\}$. In all these cases the two potentially nonidentity
foreign conditional closures are loopless. One acts only on the old
part of the base and ignores $X$; the other, belonging to $c$, only
adds $X$ from a nonempty old premise. They cannot have a disjoint
mutual-capture pair: the second closure adds no old point, so the
opposite premise has no old point; but the first fixes $\{X\}$
and cannot generate any old point from it. Looplessness then forces
both sides empty. \Cref{core:thm:two-color-feedback} excludes every
feedback family. Identity foreign closures have empty disjoint heads
and add nothing to the test.

For observer $c$ the only full-closure independent base is $B=W$,
with $Z=\{X\}$. Its two nonidentity foreign conditional closures
are precisely the pure closures of the genuine baseline $C$, so they
satisfy (FE) as well.

Finally take an identity observer. Its only full-closure base is the
whole input set, and $Z=\varnothing$. Discard empty-head queries;
they contribute nothing to any rainbow cover. In a hypothetical
remaining feedback family, a $b$-query head has at most one point of the triangle
$\{x,y,z\}$, while a $c$-query head is $\{X\}$ and its premise
contains $\{y,z,t\}$. An $a$-query without $t$ in its premise
requires at least two triangle premises. These cannot be covered by
one $b$-head and one $c$-head, as required by the rainbow condition.
Thus no $a$-head contains $t$. No other head contains $t$ either,
so no $c$-query can be present. Every remaining $a$-query has $t$
in its premise, which the remaining $b$-heads cannot cover. There
are then no $a$-queries, and no $b$-premise can be covered. This
contradicts a nonempty head union.

Additional identity colors have empty heads and cannot assist a
rainbow cover. The argument therefore proves stability under every
finite identity-owner extension.
\end{proof}
